% Dot batch 80; source SHA-256: d2862e782390adcc4088edba90732fcee8f65d44844eb667f5aea57099da61cb
% Archived source: dot_updates_20261124/batches/batch080/EGA第五卷_伪态射与可逆层延拓.tex
% Original body lines: 25--95; integration map: integration/late_integration.json
\begingroup
\def\newscope{V-L}
\def\depth{\operatorname{depth}}
\def\Spec{\operatorname{Spec}}
\def\Ass{\operatorname{Ass}}
\def\OO{\mathcal O}
\def\FF{\mathcal F}
\def\JJ{\mathcal J}
我后悔把前一命题写得既陈旧又别扭，在经典假设和自然假设之间取了个折中，却没有给出逆命题。因此我建议改作如下陈述。\cmt{L18-080-01}{法133原语较粗俗；英100编辑译注也提醒了这一点。本句是作者对行文的自评，不是撤销前一命题。}

\egas{3.3}{命题}

设 \(X\to S\) 平坦且局部有限表示，\(Q=P^\vee=\mathbf P(E^\vee)\) 是由有限型局部自由模
\(F=E^\vee\) 定义的 \(S\) 上的射影丛，记 \(P=\mathbf P(E)\)。
考虑所有由 \(Q\) 参数化的 \(X/S\) 上的除子线性系 \(D\) 所成的集合 \(\mathcal D\)，其中 \(D\) 的固定点集 \(Z\) 满足：若 \(z\in Z\)，则 \(\depth\OO_{X_s,z}\ge2\)，这里 \(s\) 是 \(z\) 在 \(S\) 中的像。
记 \(\mathcal U\) 为从 \(X\) 到 \(P\) 的相对于 \(S\) 的伪态射 \(f\) 所成的集合；要求其最大定义域 \(U=U(f)\) 满足
\[
 z\in X\setminus U\quad\Longrightarrow\quad \depth\OO_{X_s,z}\ge2,
\]
且 \(f_U:U\to P\) 满足命题3.1的非退化条件。\cmt{L18-080-02}{法133页首手写补入“第3.1中考虑的非退化条件”，正文上原“有理映射”亦手改为“伪态射”；英101已把条件写进正文。按3.1的几何纤维表述，对每个几何纤维的每个关联点，取其闭包；闭包的像不可包含在该纤维的任何超平面内。}
考虑自然映射 \(\mathcal D\to\mathcal U,\ D\mapsto f_D\)。则：

\noindent\textbf{a)} 此映射为单射。对 \(D\in\mathcal D\)，其固定点集 \(Z\) 正是 \(f_D\) 的最大定义域的补集。

\noindent\textbf{b)} 取 \(f\in\mathcal U\)，并取其一个定义开集 \(U\)，使 \(X\setminus U\) 上各点的纤维局部环深度至少为二（例如 \(U=U(f)\)）。
设 \(L_U=f_U^*\OO_P(1)\)，其中 \(f_U:U\to P\) 是 \(f\) 在 \(U\) 上诱导的态射，\(i:U\hookrightarrow X\) 是典范开浸入。
则 \(f\) 来自某个 \(D\in\mathcal D\) 的必要且充分条件，是 \(X\) 上的模 \(i_*L_U\) 可逆。\cmt{L18-080-03}{法稿排印处作缩写“il f et suffit”；英101完整译作必要且充分。这里 \(L_U\) 是拉回，不是把 \(\OO_P(1)\) 在点上求值。}

注意，若 \(X\to S\) 的纤维满足 \((S_2)\)，例如纤维正规，乃至几何正规，则上述闭子集 \(Z\) 上的深度条件只意味着：对每个 \(s\in S\)，\(Z_s\) 在 \(X_s\) 中的余维至少为二。
这时 \(\mathcal D\) 就是由 \(Q\) 参数化、无固定分支的除子线性系所成的集合。\cmt{L18-080-04}{法134末尾此处印为“由 \(P\) 参数化”，英101也随之作 \(P\)。按命题开头及关联除子的构造，参数丛始终是 \(Q=P^\vee\)，目标丛才是 \(P\)。}
另一方面，若 \(S=\Spec k\) 且 \(X\) 正规，则每个有理映射 \(f:X\dashrightarrow P\) 的定义域 \(U(f)\) 都满足
\(\operatorname{codim}(X\setminus U(f),X)\ge2\)（原稿引 II.7）。
因此在此情形，\(\mathcal U\) 由满足上述非退化条件的有理映射组成。\cmt{L18-080-05}{法134及英101末句均写“所有有理映射”，这只对深度条件成立，不能省略刚补入的非退化条件。例如 \(\mathbf A^2_k\to\mathbf P^1_k\) 的常值映射处处有定义，但其像落在一个超平面中。法136推论3.5的手写补文也把不经超平面分解写进了结论。正规情形的余维论断来自余维一局部环的赋值性质与射影目标的固有性。}

\medskip\noindent\textbf{证明。}
a) 的证明已经给出。为证 b)，先注意 \(i_*L_U\) 的形成同平坦基变换 \(S'\to S\) 交换，至少在 \(U\to X\) 拟紧时如此，原稿称可化到此情形。因此，在该情形，所需条件在忠实平坦拟紧基变换下不变。\cmt{L18-080-06}{英101保存了拟紧保留条件。法134有涉及拟紧与化归的手写边注；未辨定的字句不补造。不以任意开浸入的无条件平坦基变换为依据，见本版核查报告 L18-080。}
取 \(S'=P^\vee\)，令 \(i':U\times_S P^\vee\hookrightarrow X\times_S P^\vee\)，以撇号记拉回。
\(i'_*L'_U\) 的可逆性不因把 \(L'_U\) 换成 \(L'_U\otimes_{\OO_{S'}}M'\) 而改变，其中 \(M'\) 是 \(S'\) 上的可逆模，且张量式中省写了结构态射的拉回。这是因为
\[
 i'_*(L'_U\otimes_{\OO_{S'}} M')
 \simeq i'_*L'_U\otimes_{\OO_{S'}} M' .
\]
取 \(M'=\OO_{P^\vee}(1)\)。于是该条件也可写为 \(i'_*N'\) 可逆，其中
\[
 N'=(f_U\times_S\mathrm{id}_{P^\vee})^*
       \OO_{P\times_S P^\vee}(1,1).
\]
这里 \(\OO(1,1)\) 正是 \(P\times_S P^\vee\) 上典范关联除子 \(H\) 的可逆模，
所以 \(N'\) 恰是
\[
 D'=(f_U\times_S\mathrm{id}_{P^\vee})^*H
\]
所定义的可逆模。\cmt{L18-080-07}{法稿若干处的对偶记号及撇号不一致，英101--102进一步混淆了 \(P\) 和 \(P^\vee\)，也把除子的拉回写成无星号。此处统一按 \(X\times_S P^\vee\to P^\vee\) 及上述开浸入记号恢复。若参数丛可空，应在正秩处用其作忠实平坦覆盖；空纤维处满足深度条件的 \(X_s\) 也只能为空（否则其关联点深度为零）。}

若 \(f\) 来自 \(D\in\mathcal D\)，则 \(D'=D|_{U\times_S P^\vee}\)。
令 \(N=\OO_{X\times_S P^\vee}(D)\)，则 \(N'=N|_{U\times_S P^\vee}\)，
由上面的深度引理应用于 \(X\times_S P^\vee\to P^\vee\)，得
\(i'_*N'\simeq N\)，故 \(i'_*N'\) 可逆。\cmt{L18-080-08}{法134后改写引理号作3.2，法135又有4.2，英102亦摇摆。前一批已据上下文把直接使用的深度引理定为3.1。此处个别本子的 \(P\times_S P^\vee\to P^\vee\) 应为 \(X\times_S P^\vee\to P^\vee\)，因为 \(i'\) 及 \(N\) 在后者上。}

反过来，若满足此条件，要证明 \(f\) 来自某个 \(D\in\mathcal D\)，显然只须证明 \(D'\) 能延拓为 \(X\times_S P^\vee\) 上相对于 \(P^\vee\) 的有效 Cartier 除子。
等价地，只须证明它能延拓为该积上的有效除子 \(D\)；因为 \(U\) 含有各纤维 \(X_s\) 的所有关联点，此条件在基变换、特别是 \(P^\vee\to S\) 下保持，所以 \(D\) 自动是相对有效除子。
由上述深度引理立即可知，\(D'\) 能延拓当且仅当 \(N'\) 能延拓为可逆模，也就是 \(i'_*N'\) 可逆。

证明的末尾本应直接用充分必要条件来写，不必像我这样分两次说；此外应当先提出以下推论。

\egas{3.4}{推论（深度引理的推论）}

设 \(g:X\to S\) 平坦且局部有限表示。设 \(T\subset X\) 闭，且
\(x\in T\) 蕴含 \(\depth\OO_{X_{g(x)},x}\ge2\)。
记 \(U=X\setminus T\)，令 \(i:U\hookrightarrow X\) 为典范开浸入。
对 \(X\) 上的有限型局部自由模 \(L\)，记 \(L'=L|_U\)。

\noindent\textbf{a)} 限制函子 \(L\mapsto L'\) 全忠实，且 \(L\to i_*L'\) 为同构。
\(L\) 为秩 \(n\) 当且仅当 \(L'\) 为秩 \(n\)。

\noindent\textbf{b)} 设 \(L'\) 为 \(U\) 上的有限型局部自由模。它同构于 \(X\) 上某个有限型局部自由模 \(L\) 的限制，当且仅当 \(i_*L'\) 局部自由。\cmt{L18-080-09}{法135及英102的 b) 字面写“\(X\) 上的 \(L'\)”，但 \(i_*\) 的自变量在 \(U\) 上。由上下文承接有限型限定；即使只写局部自由，局部秩也由每个开邻域与 \(U\) 的非空交控制。}

\noindent\textbf{c)} 若 \(L'\) 是 \(U\) 上有效除子 \(D'\) 的可逆模，则 b) 的条件也是 \(D'\) 可延拓为 \(X\) 上有效除子 \(D\) 的充要条件。此延拓唯一，且为 \(D'\) 在 \(X\) 中的概形论闭包。\(D'\) 相对于 \(S\) 是有效除子，当且仅当它的延拓 \(D\) 如此。\cmt{L18-080-10}{概形论闭包由法135页下边手写补出，英102已收入正文。“当且仅当”以前面已存在延拓 \(D\) 为语境，并不主张任一相对有效除子都能延拓。}

只须应用深度引理：每个这样的 \(L\) 都满足该引理对模 \(\FF\) 所列的条件。
\par\endgroup
